\section{Viability-conditioned requirements and the limited meaning of ``ought''}

The models repeatedly generate statements of the form: if property \(V\) is to continue under disturbance class \(\D\), parameter or process \(x\) must satisfy condition \(C\).

For example,
\[
R\ge\delta V_C
\]
if a decaying capacity stock is to be maintained at \(V_C\).

Such statements can be called \emph{viability-conditioned requirements}. They have the logical form
\[
\boxed{
V\text{ is required}
\quad\land\quad
F\text{ is the relevant dynamical model}
\quad\Rightarrow\quad
\theta\in\A_V.
}
\]

No categorical moral conclusion follows. The mathematics does not select the viability criterion. It does not select the disturbance class. And it does not establish that preserving the focal organization overrides other values.

The logical order is therefore
\[
\boxed{
\begin{aligned}
\text{DESCRIPTIVE MODEL}
&\rightarrow \text{CONDITIONAL VIABILITY REQUIREMENT}\\
&\rightarrow \text{OBSERVATION}
\rightarrow \text{PRACTICAL RESPONSE}.
\end{aligned}
}
\]

This corrects a potentially misleading interpretation of the phrase ``ought before is.'' The conditional requirement can precede observing the behaviour being evaluated, but it cannot precede the dynamical knowledge from which the requirement is derived.

\subsection{Multiple requirements are constraints, not commandments}

For several viability conditions,
\[
\theta
\in
\A_1\cap\dots\cap\A_k.
\]
There is no mathematical reason the intersection must be non-empty. Consequently, viability theory does not generate a universal harmony among system requirements. Sometimes the architecture itself must change.

\section{Why this is not yet a theory of morality}

The maintenance language has tempting human translations. Reliable information resembles honesty. Return-loop closure resembles reciprocity. Negative feedback on damaging behaviour resembles accountability. Preserving a valuable relation after recoverable failure resembles forgiveness or repair. Maintaining model revisability resembles corrigibility.

These analogies may eventually be useful. But they are not results of the present paper.

Morality-as-Cooperation already provides a developed account of moral rules as biological and cultural solutions to recurrent cooperation problems and has generated cross-cultural predictions about such rules \citep{curry2019mapping,curry2019societies}.

A distributed-maintenance account would need to do more than relabel established cooperative behaviours. It would need to generate discriminating predictions---for example, conditions under which a norm preserves collective learning or correction despite not directly solving a standard cooperation game.

Until such predictions exist, human ethical mappings should be treated as hypotheses rather than evidence for the framework.

\section{An empirical programme}

A general language becomes scientifically useful only if it changes what is measured or predicted. We therefore propose an evidence hierarchy.

\subsection{Level 1: analogical mapping}

A researcher can describe a system using maintenance vocabulary. This has minimal evidential value.

\subsection{Level 2: independently specified viability boundary}

The relevant continuation condition is defined independently of the outcome used to test the model. Examples might include an engineering tolerance, a transmission threshold, an extinction boundary, or a protocol soundness requirement.

\subsection{Level 3: measured corrective capacity}

Systems with similar present state but different capacity should have different responses to controlled or naturally occurring disturbances. This directly tests
\[
\text{state}\neq\text{capacity}.
\]

\subsection{Level 4: identified renewal pathway}

The processes replenishing corrective capacity are independently measured. Intervening on the pathway should alter later disturbance tolerance in the predicted direction.

\subsection{Level 5: function--proxy discrimination}

Where renewal is allocated using a proxy \(Y\), the framework predicts failure modes when the relationship between \(Y\) and function \(F\) degrades. A strong test would manipulate that relationship while holding other conditions as constant as possible.

\subsection{Level 6: prospective improvement over state-only models}

The strongest practical test is comparative. Let \(\mathcal M_0\) use present state only, \(\mathcal M_1\) use present state and empirical trend, and \(\mathcal M_2\) add explicit corrective capacity, renewal and function--proxy variables.

The framework gains empirical support only if \(\mathcal M_2\) improves prospective prediction or intervention beyond simpler alternatives. Otherwise the additional maintenance language has not demonstrated explanatory value.

\section{Testable hypotheses generated by the synthesis}

The framework suggests several hypotheses that can fail.

\paragraph{H1. State-matched capacity divergence.}
For systems matched on current performance,
\[
X_1\approx X_2,
\]
independently measured corrective capacity should predict different responses to disturbance.

\paragraph{H2. Renewal-path intervention.}
Reducing a genuine renewal pathway should reduce future corrective capacity before necessarily altering current performance. Restoring the pathway should reverse that effect if irreversible damage has not occurred.

\paragraph{H3. Failure-domain diversity.}
Nominal redundancy should predict resilience less well than redundancy weighted by independent failure domains when correlated failure is material.

\paragraph{H4. Proxy-gap instability.}
Where reinforcement is driven by proxy \(Y\), deterioration in the relationship
\[
Y\leftrightarrow F
\]
should predict increasing allocation to processes with declining viability contribution.

\paragraph{H5. Mixed-sign interventions.}
Interventions that improve one viability margin while worsening another should resist stable ranking unless an explicit weighting or higher-level constraint is supplied.

\paragraph{H6. Architectural infeasibility.}
When admissible regions for necessary viability loops cease to intersect, parameter optimization within the current architecture should systematically fail, whereas an architectural change can restore feasibility.

These hypotheses are considerably stronger than the statement that all persistent systems require maintenance.
